El trabajo partía de un problema concreto: el alto coste de mantener
sincronizados backend, comunicación e interfaz ante un modelo de datos en
evolución. La solución construida demuestra que ese coste puede reducirse de
forma drástica si el modelo ORM se convierte en la única fuente de verdad y el
resto del sistema se deriva de él mediante un contrato uniforme y una cadena
de generación automática.

De los cinco objetivos específicos planteados, cuatro se han completado y
consolidado con pruebas (API modular, modelado de datos, comunicación
estandarizada y generación dinámica de formularios) y el quinto --- generación
automática de esquemas desde el ORM --- se ha alcanzado parcialmente, con su
base funcional y su versionado operativos pero con extensiones pendientes. El
grado de consecución de los objetivos generales es, por tanto, alto: existe un
framework Backend/Frontend reutilizable, con una arquitectura desacoplada del
dominio y adaptada de forma evolutiva a las necesidades reales del \gls{itc} y
del \gls{jbvc}. La metodología iterativa incremental resultó especialmente
adecuada, pues su planificación evolutiva absorbió la incertidumbre inicial
del backend y permitió cerrar cada incremento con una versión funcional y
verificada.
